\documentclass{article}
\usepackage[utf8]{inputenc}
\usepackage{amsmath}
\usepackage{amssymb}

\begin{document}
\title{Interdiction of cyber attacks}
\maketitle
\section{Sets and Parameters}
\begin{itemize}
    \item $N$: Set of nodes;
    \item $L$: number of levels $- 1$;
    \item $N^I \subseteq N$: Set of nodes that can be breached from the outside (level 1);
    \item $N^R \subseteq N$: Set of nodes that bring a revenue if reached (level $L$);
    \item $r \in N$; root node (level 0) , it is connected to all nodes in $N^I$, and only to them:
    \item $lev(i)$: Level of node $i \in N$;
    \item $A$: Set of available arcs;
    \item $c_{ij}$: Cost for attacking arc $(i, j) \in A$;
    \item $d_{ij}$: Cost for defending arc $(i, j) \in A$;
    \item $B_a$: Attacker's total budget;
    \item $B_d$: Defender's total budget;
    \item $r_j$: Reward associated with node $j \in N^R$.
\end{itemize}


\section{Defender Outer Problem}
The defender seeks to minimize the attacker's maximum profit by interdicting arcs, subject to a budget $B_{d}$.

\subsection{Variables}
\begin{itemize}
    \item $y_{ij}$: 1 if the defender interdicts arc $(i,j) \in A$, 0 otherwise;
    \item $w_{ijk}$: 1 if the defender interdicts arcs $(i,j), (j,k) \in A$ at the same time, 0 otherwise;
    \item $z$: continuous variable containing the maximum profit collectable by the attacker given the current interdiction strategy ($y$ variables).
\end{itemize}

\subsection{Model}\label{mo}
\begin{align}
    (O) \  \min &\ z\label{c0} \\
    \text{s.t.} & \sum_{(i,j) \in A} d_{ij} y_{ij} \leq B_{d} \label{c1} \\
    & w_{ijk} \geq y_{ij} + y_{jk} - 1 & \forall (i,j), (j,k) \in A \label{c2a} \\
    & w_{ijk} \le y_{ij}  & \forall (i,j), (j,k) \in A \label{c2b} \\
    & w_{ijk} \le y_{jk} & \forall (i,j), (j,k) \in A \label{c2c} \\
    & z \geq v^p - \sum_{(i,j) \in A} u_{ij}^p y_{ij} + \sum_{(i,j,), (j,k) \in A} u_{jk}^p w_{ijk} & \forall p \in \mathcal{P} \label{c3} \\
    & y_{ij} \in \{0, 1\} & \forall (i,j) \in A \label{c4} \\
    &  w_{ijk} \in \{0, 1\} & \forall (i,j), (j,k) \in A\label{c5} \\
    & z \geq 0  \label{c6} 
\end{align}
The objective function (\ref{c0}) minimizes the value of the technical variable $z$.
The inequality (\ref{c1}) ensures that the interdiction budget of the defender is not exceed.
Constraints  (\ref{c2a})-(\ref{c2c}) keep the value of $w$ variables aligned with those of $y$ variables.
Inequalities (\ref{c3}) implement the assignment of the maximum profit collectable by the attacker given the current interdiction plan to variable $z$. The set $\mathcal{P}$ contains all possible strategies of the attacker. The constant $v^p$ contains the total profit collected if the attacker strategy $p$ is implemented and no interdiction is applied, while the constant $u^p_{ij}$ contains the potential profit carried over arc $(i,j)$ in the attack plan $p$. Notice that the term $u_{ij}^p y_{ij}$ represents the reduction in the attacker's profit when the defender chooses to interdict arc $(i,j)$, while the term $u^p_{jk} w_{ijk}$ is a correction factor (recuperation) used when two consecutive arcs $(i,j)$ and $(j,k)$ are interdicted, preventing over-counting of the interdiction effect in specific path structures.
Constraints (\ref{c4})-(\ref{c6}) define the domain of the variables.

Since there is potentially a huge number of constraints (\ref{c3}), one for each possible plan of the attacker, it is not feasible to solve the model as it is. For this reason, only the relevant constraints (\ref{c3}) will be generated iteratively as Benders' cuts.

Notice that each iteration the optimal solution value of model $O$ cannot be lower than the previous iteration. At the same time, solving each of the inner problem (see Section \ref{in}) will produce a valid upper bounds for the optimal solution of $O$. Once upper and lower bounds meet, the iterative algorithm can be stopped and the optimal solution has been found.

\section{Attacker Inner Problem}\label{in}

We consider an interdiction plan $y^p$ of the defender at iteration $p$, obtained by solving the model presented in Section \ref{mo} with a subset of constraints (\ref{c3}), such that  $y^p=1$ if arc $(i,j) \in A$ is interdicted, $0$ otherwise. A model that produces the plan with highest profit for the attacker is described. 

Formally, we define the following set of arcs that are not interdicted by the defender: $F^p=  \{(i,j) \in A: y^p_{ij}=0\}$, which will be used to devise an optimal attack plan that will exploit the arcs available in $F^p$ only, to respond to the defence plan. 

\subsection{Decision Variables}
\begin{itemize}
    \item $x_{ij}$: 1 if the arc $(i,j) \in A$ is part of the attack plan, 0 otherwise;
    \item $u_{ij}$: continuous variable containing the profit carried by art $(i,j) \in A$ in the current attack plan.
\end{itemize}

\subsection{Model}

\begin{align}
    (I^p) \ v= \max &\ \sum_{(r,j) \in F^p} u_{rj} \label{d0}  \\
    \text{s.t.\ } & \sum_{(i,j) \in F^p} c_{ij} x_{ij} \leq B_a \label{d1} \\
    & \sum_{(j,i) \in F^p} x_{ji} \leq 1 & \forall i \in N  \backslash \{r\} \label{d2} \\
    & u_{ij} \leq \sum_{(j,k) \in F^p} u_{jk} & \forall j \in N \backslash N^R \backslash \{r\}, \forall (i,j) \in F^p \label{d3} \\
    & u_{ij} \leq M x_{ij} & \forall j \notin N^R,  \forall (i,j) \in F^p \label{d4} \\
    & u_{ij} \leq r_j x_{ij} & \forall j \in N^R, \forall (i,j) \in F^p \label{d5} \\
    & x_{ij}  \in \{0,1\} & \forall (i,j) \in F^p \label{d6} \\
    & u_{ij} \ge 0 & \forall (i,j) \in F^p \label{d7}
\end{align}

The objective function (\ref{d0}) maximizes the total revenue collected by the attack plan, formally the total revenue at the root node $r$..
The inequality (\ref{d1}) ensures that the interdiction budget of the attacker is not exceed.
Constraints  (\ref{d2}) forces the attack plan to be an arborescence, by forcing that at most one active arc can enter each node $i \in N \backslash \{r\}$.
Inequalities (\ref{d3}) regulate the collection of the profit along the active arborescence.
Inequalities (\ref{d4}) and (\ref{d5}) impose that profit can only be collected on arcs associated with active arcs. The first set regulates arcs not involving leaves of the arborescence, while the second set refers to arcs entering leave nodes. The constant $M$ is a big number that in the case of the problem under investigation can be set as $M=\sum{i \in N^R} r_i$.
Constraints (\ref{d6}) and (\ref{d7}) define the domain of the variables.

Once model $I^p$ is solved, the optimal values of the relevant variables $u^p$, together with the optimal value $v^p$, are used to build a new inequality of type (\ref{c3}) that is added to model $O$.

\subsubsection{Valid Inequalities}
The following set of valid inequality can be added to model $I^p$ to speed up computation:
\begin{equation}
\sum_{(i,j) \in A: level(i)=l \land level(j)=l+1} u_{ij}=v \ \ \ \forall l \in \{0, 1, \dots, L-1\} \label{extra}
\end{equation}
Equalities (\ref{extra}) impose that the profit passed from one level of the graph to the level on top of it has to equal $v$.
\end{document}
